\abstractch

计算机网络课程材料包含大量非结构化知识，人工手动整理结构化文本耗时长，难以满足课程教学与研究的需要。课程文档以非结构化 PDF（Portable Document Format）与 PPT（PowerPoint）文档为主，包含大量专业术语与跨章节的复杂语义关系，不利于构建完整准确的知识图谱。实现课程知识的自动化抽取与结构化表示，对于提升教学资源的智能化利用水平具有重要意义。

为此，设计并实现一套面向计算机网络课程材料的自动化知识图谱构建系统十分必要。以大型语言模型为核心，通过提示工程方法完成课程文本的命名实体识别与关系抽取，将非结构化文档转化为结构化知识图谱，为课程教学提供可查询的知识基础。

七阶段流水线实现了从原始文档到知识图谱的端到端转换。文档预处理阶段使用多模态视觉语言模型将课程文档转换为结构化 Markdown。长文档处理采用标题层级的递归分块算法，结合词元预算控制，在保证语义连贯的前提下实现文本切分。知识抽取由实体识别和关系抽取两个子任务协同完成，分别抽取四类实体与四类关系。实体歧义通过稠密向量嵌入的规范化方法解决：利用 FAISS 向量索引与并查集算法完成实体同一性判定与等价类归并。知识融合后的三元组导入 Neo4j 图数据库，通过交互式前端界面提供关系导航与语义查询功能。

实验基于 205 份计算机网络课程 PDF 文档。原始文档经 VLM 转换为 Markdown 后，分块生成 16,438 个段落级文本块，包含多级标题内容，总计约 1,040 万词元作为知识抽取的结构化输入。在命名实体识别与关系抽取层面，启用LLM复核的最终知识图谱（\texttt{merged-fs-recheck}）包含 32,760 个实体（概念 22,696 个、机制 6,406 个、性能指标 2,192 个、协议 1,466 个）和 17,627 条关系三元组。在500条独立金标上采用"严格／软匹配／槽位／类型级"四档匹配口径评估，软匹配三元组F1为15\%~19\%，槽位F1为37\%~42\%，类型级F1可达74\%~82\%。

实验结果表明，所构建的自动化流水线能够从大规模非结构化课程材料中有效提取结构化知识，构建层次清晰、覆盖面广的知识图谱，验证了大型语言模型在课程知识建模与抽取中的可行性与有效性。

\vspace{\baselineskip}

\keywordsch 大语言模型；命名实体识别；关系抽取；知识图谱；课程知识；提示工程；实体规范化
